Phần còn lại của báo cáo đồ án tốt nghiệp được tổ chức như sau.

Chương 2 trình bày nền tảng lý thuyết và các nghiên cứu liên quan đến quan sát lớp học, nhận diện hành vi bằng thị giác máy tính, khám phá liên kết thời gian trên chuỗi thời gian và vai trò của mô hình ngôn ngữ lớn trong hỗ trợ khám phá đồ thị. Chương này cũng tổng hợp khoảng trống nghiên cứu làm cơ sở cho thiết kế CausClass.

Chương 3 mô tả phương pháp đề xuất, bao gồm tổng quan quy trình, kiến trúc YOLO26n-MHSA cho nhận diện hành vi, bước theo dõi và xây dựng chuỗi thời gian, cách trích xuất đồ thị nền bằng AERCA và thủ tục beam search có hướng dẫn bởi LLM.

Chương 4 trình bày thiết kế thực nghiệm và kết quả đánh giá. Nội dung chương bao gồm dữ liệu cho Stage 1, phân tích ablation cho bộ phát hiện hành vi, dữ liệu cho Stage 2, phân tích ablation cho khám phá đồ thị và một nghiên cứu tình huống trên video lớp học thật.

Chương 5 tổng kết các kết quả chính của đồ án, nêu rõ những giới hạn hiện tại và đề xuất các hướng phát triển tiếp theo cho CausClass trong bối cảnh phân tích lớp học dựa trên bằng chứng.